\introduction

Развитие робототехники и искусственного интеллекта в последние десятилетия привело к широкому распространению мульти-роботизированных систем (МРС) в различных областях деятельности человека. Группы роботов, действующих совместно, применяются в складской логистике, автономном транспорте, поисково-спасательных операциях, сельском хозяйстве и многих других сферах~\cite{wurman2008kiva, parker2008mrs}. Использование нескольких автономных агентов вместо одного позволяет повысить производительность за счёт параллельного выполнения задач, обеспечить отказоустойчивость системы и масштабируемость решений. Координация множества автономных агентов, работающих в общем пространстве, представляет собой одну из ключевых задач современной робототехники.

Управление мульти-роботизированными системами включает два фундаментальных класса задач: распределение задач между роботами (Multi-Robot Task Assignment, MRTA) и поиск бесконфликтных маршрутов для группы агентов (Multi-Agent Path Finding, MAPF). Задача MRTA определяет, какой робот выполняет какую задачу, а задача MAPF обеспечивает построение маршрутов движения роботов к назначенным целям таким образом, чтобы агенты не сталкивались друг с другом. Для решения каждой из этих задач разработано значительное число алгоритмов, существенно различающихся по архитектуре, вычислительной сложности, оптимальности получаемых решений и области применимости~\cite{stern2019mapf, gerkey2004taxonomy}.

\textbf{Актуальность} работы обусловлена растущей потребностью в инструментах для объективного сравнения алгоритмов управления МРС. Существующие реализации алгоритмов MAPF и MRTA, как правило, разрозненны, выполнены в различных программных средах и на разных языках программирования, что затрудняет их прямое сравнение в единых условиях. Кроме того, большинство доступных инструментов ориентированы на исследователей и не предоставляют интуитивного графического интерфейса для настройки экспериментов, визуализации работы алгоритмов и автоматизированного сбора метрик. Выбор подходящего алгоритма для конкретной прикладной задачи остаётся нетривиальной проблемой, поскольку эффективность каждого алгоритма существенно зависит от параметров среды: размера карты, количества роботов, плотности препятствий и характера выполняемых задач.

\textbf{Целью} данной работы является создание кроссплатформенного программного комплекса для тестирования и сравнения эффективности алгоритмов управления мульти-роботизированными системами.

Для достижения поставленной цели необходимо решить следующие \textbf{задачи}:

\begin{enumerate}[arabic]
    \item провести анализ существующих алгоритмов решения задач MAPF и MRTA, их классификацию и особенности применения;
    \item разработать архитектуру программного комплекса, обеспечивающую модульность и расширяемость;
    \item реализовать алгоритмы поиска бесконфликтных маршрутов: CBS, CBS-SIPP, ECBS, EECBS, M*, CA*, WHCA*, PIBT;
    \item реализовать алгоритмы распределения задач: венгерский алгоритм, жадный алгоритм, последовательный аукцион, минимальная стоимость потока, комбинаторный аукцион;
    \item разработать графический интерфейс пользователя, включающий редактор карт, панель выбора алгоритмов и визуализацию симуляции;
    \item реализовать режим пакетного тестирования для автоматизированного сравнительного анализа алгоритмов при различных параметрах среды;
    \item обеспечить кроссплатформенность программного комплекса;
    \item провести сравнительный анализ реализованных алгоритмов по ключевым метрикам эффективности.
\end{enumerate}

\textbf{Объектом исследования} является программный комплекс для оценки эффективности алгоритмов управления мульти-роботизированными системами. \textbf{Предметом исследования} являются алгоритмы управления мульти-роботизированными системами и их сравнительная эффективность при различных конфигурациях среды и количестве агентов.

\textbf{Результатом работы} является кроссплатформенный программный комплекс на языке Python с графическим интерфейсом на базе библиотеки Pygame, реализующий восемь MAPF-алгоритмов и пять MRTA-алгоритмов. Комплекс предоставляет интерактивную среду для создания сценариев на сеточных картах, визуализации работы алгоритмов в режиме реального времени и проведения автоматизированного пакетного сравнительного анализа с генерацией отчётов и графиков по ключевым метрикам: времени выполнения (makespan), суммарной стоимости путей (SOC), времени планирования и потреблению памяти. С помощью PyInstaller обеспечена сборка в автономный исполняемый файл для операционных систем Windows, Linux и macOS.
